\section{数据工程 I：可验证任务与不可验证任务}

讲者将数据分为两大类：\textbf{verifiable} 与 \textbf{non-verifiable}。前者如数学、代码执行类任务，答案可以通过程序、规则或形式系统验证；后者如写作、开放问答、安全边界判断，往往需要 rubric 或模型评估。

这一划分直接决定训练策略。对于可验证任务，RL 可以利用高精度自动反馈并开展规模化 rollout；对于不可验证任务，核心是构建稳定的 rubric 和评估一致性，避免奖励噪声吞噬学习信号。

\begin{knowledgebox}{两类数据的本质差异}
\textbf{Verifiable}：反馈精确、可并行、大规模自动化，适合高频迭代。\\
\textbf{Non-verifiable}：反馈主观、漂移风险高，需要多 grader 交叉与人工抽检。
\end{knowledgebox}

讲者特别提到安全相关任务中 ``对某些用户安全、对另一些用户不安全'' 的差异，意味着这类任务很难用单一二元标签定义奖励。工业做法通常是把高风险维度拆成多个约束项并进行分层扣分。

\begin{warningbox}{把所有任务都强行转成可验证会带来偏差}
过度追求可验证会导致模型偏向 ``可测任务''，忽略真实用户需求中大量灰度场景。结果是在离线 benchmark 提升、线上体验下降。
\end{warningbox}

\begin{importantbox}{数据定义优先于算法细节}
讲者多次强调 ``data is the most important''。在同等算力下，数据分层和数据路由策略往往比更换 RL 目标函数更快带来稳定收益。
\end{importantbox}

从课程语境看，数据工程不是后处理环节，而是训练主干。尤其在 agent 场景，数据质量决定了模型是否能学会 ``行动责任''，而不只是 ``文本正确''。

\subsection{本章小结}
可验证与不可验证任务需要不同训练机制。高质量 agent 训练首先是数据问题，其次才是优化问题。

